\scanpage{41}\mspage{40}
\deleted{\inserted{car $C$ est $ {\unsure{\tau|\mathcal T}}'$-fermé}}
\[
\begin{aligned}
 &\deleted{C=}\ \underset{\substack{\text{\deleted{fermé de}}\\\deleted{(T,\tau'_0)}\\\text{\deleted{soit }}\deleted{B}}}{\underset{\rule{.35pt}{.35em}\rule{4.4em}{.35pt}\rule{.35pt}{.35em}}{\deleted{(C\cap T)}}}
 \ \deleted{\cup}\ 
 \underset{\substack{\text{\deleted{\unsure{fermé de|partie de}} \deleted{\inserted{partie de $S$ dont la trace sur}}}\\
 \text{\deleted{\inserted{$T$ est $\unsure{\tau'_0|\tau_0}$-fermée, et}}}\\
 \text{\deleted{$(S,\tau)$, soit} \deleted{\inserted{la trace sur $S\setminus T$}}}\\
 \text{\deleted{\inserted{est égale à $S\setminus T$}}}\\\deleted{A}}}{\underset{\rule{.35pt}{.35em}\rule{7.8em}{.35pt}\rule{.35pt}{.35em}}{\deleted{\overline{(C\cap(S\setminus T))}^{(\tau')}}}}
\end{aligned}
\]
\deleted{OK. Je voudrais résumer ce qui a été vu dans un lemme topologique plus général :}

\uline{Lemme}\quad \unsure{Soit|Soient} $(\deleted{\unsure{T}}S,\tau)$ espace topologique, $T\deleted{\unsure{_0|_{\varnothing}}}\subset S$ une partie fermée, et $\tau'_0$ une topologie \deleted{\unsure{p}} sur $T$ plus fine que la top. induite (i.e.
\[
 (T,\tau'_0)\mathbin{\unsure{-|\longrightarrow}}(S,\tau)\qquad\text{continue).}
\]

\deleted{Soit $\tau'$ une topologie sur}

\deleted{Il y a alors \unsure{d'}} Soit \uline{$\unsure{\mathscr{T}|\mathcal T|\tau}$} l'ens des topologies sur $S$ qui sont

$1^\circ)$ plus fines que $\tau$, et

$2^\circ)$ induisent $\tau'_0$ sur $T$.

\textcircled{a} Il y a alors dans \uline{$\mathscr T$} une topologie la moins fine que toutes, soit $\tau'_{\unsure{a|\ast}}$.

\textcircled{b} Pour qu'une partie $X$ de $S$ soit fermée pour $\tau$, il f et s. que (posant $U=S\setminus T$)
\[
\begin{aligned}
1^\circ)\quad &\deleted{\unsure{A}}X\cap \deleted{\unsure{B}}U\deleted{\unsure{\longrightarrow}}\overset{\inserted{\mathrm{dfn}}}{=}X_{\unsure{u|U}}\quad\text{soit fermé dans }U\quad(\text{pour }\tau\mid U)\\
2^\circ)\quad &X\deleted{\unsure{_a}}\cap\deleted{\unsure{B}}T\overset{\inserted{\mathrm{dfn}}}{=}X_{\deleted{\unsure{B}}T}\ \text{soit fermé dans }T\quad(\text{pour }\tau'_0)\\
3^\circ)\quad &\deleted{\overline{\unsure{X\cap U}}}\ \overline{X_u}^{(\tau)}\deleted{\unsure{\subset}}{}\cap T\subset X_{\deleted{\unsure{B}}T}.
\end{aligned}
\]
\deleted{\textcircled{c} Les fermés de $\tau'$ \unsure{sont|sont alors} \unsure{écrits} de façon canonique sous la forme \unsure{$+$|$t$}}
